% Page 016

\seite{16}

\margmark{Maibaum}%
                16

\pend
\jahresueberschrift{Das Jahr 1888}
\pstart

Das Jahr 1888 wies die Schülerzahl 43 auf.\ob
Am 9.\ März, am Todestage Kaiser Wilhelms I. fand\ob
eine ansprechende Todes-Feier Statt; ebenso am 16.\ Juni\ob
am Tage nach dem Todestage des alten Kaisers\ob
Friedrichs III.

In dem Gelände von Sefferweich wurden im letzten\ob
Jahre die Forstpflanzungen abgehalten. Die Gegend und\ob
Hutflächen die Pflanzfläche scheint für diese Anlage\ob
ganz besonders geeignet zu sein.

Noch ist ein alter Usus nicht unerwähnt gelassen: bei Gelegenheit\ob
des\unsicher, daß die erwachsenen Jungen den Mädchen am\ob
1.\ Mai einen Maibaum errichten. Dazu mußten sie den\ob
\margmark{beim Maibaum\\1851.}%
schönsten eine Buche. Es wurde mit Bändern und Schleifen\ob
geziert und mitten im Dorfe aufgerichtet. Als Gegen\ohb
gabe spendeten die Mädchen ein Fest, bestehend in\ob
Kaffee mit Kuchen und hinterdrein Schwarzkirschwein.\ob
Im Jahre 1851 wurde der hübsch geschmückte des Nachts\ob
gestohlen. Den folgenden Tag, an einem Sonntag,\ob
fanden sie den Baum als umgehauen. Infolge dessen\ob
wurde über die Täter verhandelt und sollte den\ob
Thäter für 3 Mark Strafgelder zurückgegeben werden. Bei\ob
der Zurückgabe spielte der damalige Kirchendiener kennen\ob
solle; die anwesenden Gäste gerieten in Tumult\ob
und eine große Schlägerei auf dem Saale war die Folge.\ob
Zugestellten sonderten 14 Personen, verurteilt und sollten die\ob
gerichtlichen Kosten (ca. 500) zu tragen. Seitdem war\ob
den Spaß am Maibaum verloren.
\pend
